% !PRINTED 501
\weberpage{501}

If in this representation numerator and denominator are imprimitive, one can take out the divisor
and obtains a representation of the form
\weq{2}{A = a\,\frac{E_1(x,y,z,\ldots)}{E_2(x,y,z,\ldots)} = a\,\frac{E_1}{E_2} = aE,}
in which $a$ is a \gesperrt{positive}, integral or fractional rational number,
while $E_1, E_2$ are primitive functions in $R$. From this
it follows:

1. \gesperrt{Every rational functional can be transformed,
by multiplication with a primitive function in $R$, into an
integral function in $R$, and several rational functionals
can be represented as fractions whose common denominator
is a primitive function.}

The positive number $a$ and the quotient $E_1:E_2$ are completely
determined by $A$. For if we put $a$ into the form of a
rational fraction $q_1:q_2$ and assume that
\[
\frac{q_1}{q_2}\,\frac{E_1}{E_2} = \frac{q_1'}{q_2'}\,\frac{E_1'}{E_2'},
\]
then it follows that
\[
q_1 q_2' E_1 E_2' = q_2 q_1' E_2 E_1'.
\]
% [sic: Weber prints ``Functionen $R$'' without ``in''.]
Here, then, we have two integral functions $R$ with integral
coefficients which are equal to one another, and whose divisor,
since $E_1 E_2'$ and $E_2 E_1'$ are primitive functions, is $q_1 q_2'$
or $q_2 q_1'$. Consequently $q_1 q_2' = q_2 q_1'$, and hence
also
\[
\frac{q_1}{q_2} = \frac{q_1'}{q_2'}, \quad \frac{E_1}{E_2} = \frac{E_1'}{E_2'}.
\]

2. \gesperrt{We call the positive rational number $a$ the
absolute value of the functional $A$.}

In the case in which $A$ is itself a number, $a$ is the absolute
value of $A$ in the ordinary sense of this word, and $E$
is to be set $= +1$ or $= -1$, according as $A$ is positive or
negative. Thus $E$ is nothing more than the sign of
$A$. In the far-reaching generalisation of these elementary
concepts lies the strangeness which our definition offers at first
glance. It will, however, prove in what follows to be thoroughly
appropriate and useful.

From \S.\,135, 3. there results the theorem:
%CONT>>
